\section*{Chapter \ref{ch:b1:income-and-substitution-effects} --- Income and Substitution Effects}
% Numbers recomputed in code/bachelor-1/06-income-and-substitution-effects/exercises/solutions.py (tested).

\begin{solution}{exo:b1:income-and-substitution-effects:slutsky-cd}
Before: $x_1 = 100/2 = 50$. Slutsky compensation: $100 + (2 - 1) \times 50 = 150$.
At prices $(2, 1)$ with income 150: $x_1 = 150/(2 \times 2) = 37.5$. Without
compensation: $x_1 = 100/4 = 25$. \omterm{def:b1:income-and-substitution-effects:effects}{Substitution effect} $37.5 - 50 = -12.5$;
\omterm{def:b1:income-and-substitution-effects:effects}{income effect} $25 - 37.5 = -12.5$; total $-25$.
\end{solution}

\begin{solution}{exo:b1:income-and-substitution-effects:inferior}
The \omterm{def:b1:income-and-substitution-effects:effects}{substitution effect} of a price rise is negative or zero for every good
(\cref{prop:b1:income-and-substitution-effects:negative}). The price rise
lowers what income can buy, and for an \omterm{def:b1:consumer-choice:normal-inferior}{inferior good} less real income means
more of the good: the \omterm{def:b1:income-and-substitution-effects:effects}{income effect} is positive. The total is zero if the
two effects cancel exactly, the boundary between an ordinary and a \omterm{def:b1:income-and-substitution-effects:giffen}{Giffen
good}.
\end{solution}

\begin{solution}{exo:b1:income-and-substitution-effects:read}
Sweden: $1.554/2.149 = 0.723$, a fall of $27.7\%$. The ten countries:
$1.928/1.926 = 1.001$, a rise of $0.1\%$. The difference is not the effect
of the tax: biofuels, fuel efficiency, the mix of cars and incomes changed
differently in Sweden and in the other countries over thirty years, and the
mean of ten countries was not built to track Sweden before 1991
(\cref{exo:b1:income-and-substitution-effects:chart}).
\end{solution}

\begin{solution}{exo:b1:income-and-substitution-effects:expenditure-cd}
Minimise $p_1x_1 + p_2x_2$ subject to $x_1x_2 = u$. The tangency
$x_2/x_1 = p_1/p_2$ gives $x_2 = p_1x_1/p_2$, so $p_1x_1^2/p_2 = u$:
$h_1 = \sqrt{u\,p_2/p_1}$ and $h_2 = \sqrt{u\,p_1/p_2}$. Then
$e = p_1h_1 + p_2h_2 = 2\sqrt{u\,p_1p_2}$. With $v(p, m) = (m/2p_1)(m/2p_2) = m^2/(4p_1p_2)$,
$e\bigl(p, v(p, m)\bigr) = 2\sqrt{m^2/4} = m$.
\end{solution}

\begin{solution}{exo:b1:income-and-substitution-effects:hicks}
Utility before is $50 \times 50 = 2500$; at prices $(2, 1)$ it costs
$e = 2\sqrt{2500 \times 2} = 141.42$. Then $x_1 = 141.42/4 = 35.36$. \omterm{def:b1:income-and-substitution-effects:effects}{Substitution
effect} $35.36 - 50 = -14.64$; \omterm{def:b1:income-and-substitution-effects:effects}{income effect} $25 - 35.36 = -10.36$. The Slutsky
compensation, 150, is larger than the Hicks one, 141.42, and good 1 is
normal, so the Slutsky consumer buys more good 1 (37.5) and her
\omterm{def:b1:income-and-substitution-effects:effects}{substitution effect} is smaller in size.
\end{solution}

\begin{solution}{exo:b1:income-and-substitution-effects:shephard}
$\partial e/\partial p_2 = 2\sqrt{u\,p_1} \times \tfrac{1}{2\sqrt{p_2}} = \sqrt{u\,p_1/p_2} = h_2$.
$\partial h_1/\partial p_2 = \tfrac12\sqrt{u/(p_1p_2)}$ and
$\partial h_2/\partial p_1 = \tfrac12\sqrt{u/(p_1p_2)}$: equal, and positive, so the
two goods are \omterm{def:b1:income-and-substitution-effects:net}{net substitutes}.
\end{solution}

\begin{solution}{exo:b1:income-and-substitution-effects:fuel}
By the elasticity form of the \omterm{thm:b1:income-and-substitution-effects:slutsky}{Slutsky equation},
$\varepsilon^c = \varepsilon + s\eta = -0.35 + 0.06 \times 0.8 = -0.35 + 0.048 = -0.302$.
A compensated $20\%$ rise lowers fuel use by about $0.302 \times 20\% = 6.0\%$,
against $0.35 \times 20\% = 7.0\%$ without the cheque.
\end{solution}

\begin{solution}{exo:b1:income-and-substitution-effects:cournot}
$x_1 = am/p_1$ and $x_2 = (1 - a)m/p_2$: $\varepsilon_{11} = -1$, $\varepsilon_{21} = 0$,
$\eta_1 = \eta_2 = 1$; the shares are $s_1 = a$, $s_2 = 1 - a$. Engel:
$a \times 1 + (1 - a) \times 1 = 1$. Cournot for $j = 1$:
$a \times (-1) + (1 - a) \times 0 = -a = -s_1$.
\end{solution}

\begin{solution}{exo:b1:income-and-substitution-effects:chart}
Three reasons: road emissions per person also moved with fuel efficiency,
biofuel blending and the mix of petrol and diesel cars; Sweden had a deep
recession in the early 1990s that cut driving for reasons of income; the
other countries had fuel taxes and policies of their own, and differ from
Sweden in ways that change over time. A comparison must build a group, or a
weighted mix of countries, whose emissions followed Sweden's closely
\emph{before} the tax, so that after it the gap can be read as the effect
--- the synthetic control of the record box --- and it must check that the
gap is not produced by any one country or by chance.
\end{solution}

\begin{solution}{exo:b1:income-and-substitution-effects:elasticity}
Multiply the \omterm{thm:b1:income-and-substitution-effects:slutsky}{Slutsky equation} by $p_j/x_i$, and use $h_i = x_i$ at the point:
\[
\frac{p_j}{x_i}\frac{\partial x_i}{\partial p_j} = \frac{p_j}{h_i}\frac{\partial h_i}{\partial p_j} - \frac{p_jx_j}{m}\cdot\frac{m}{x_i}\frac{\partial x_i}{\partial m},
\]
that is, $\varepsilon_{ij} = \varepsilon^{c}_{ij} - s_j\eta_i$.
\end{solution}

\begin{solution}{exo:b1:income-and-substitution-effects:giffen-share}
A \omterm{def:b1:income-and-substitution-effects:giffen}{Giffen good} has $\varepsilon_{ii} > 0$, so
$\varepsilon^{c}_{ii} - s_i\eta_i > 0$, that is $-s_i\eta_i > -\varepsilon^{c}_{ii} \ge 0$:
$\eta_i < 0$ (inferior) and $s_i|\eta_i| > |\varepsilon^{c}_{ii}|$. \omterm{prop:b1:income-and-substitution-effects:aggregation}{Engel aggregation}
then gives $\sum_{j \neq i} s_j\eta_j = 1 - s_i\eta_i = 1 + s_i|\eta_i| > 1$. With a
small share $s_i$, the condition needs a very large $|\eta_i|$ or an almost
zero \omterm{def:b1:income-and-substitution-effects:effects}{substitution effect}; staples that take most of a poor household's
budget are where it can be met.
\end{solution}

\begin{solution}{exo:b1:income-and-substitution-effects:hunan}
The vouchers lowered the price of rice. Giffen behaviour means that rice
consumption falls when its price falls. A rise in rice consumption among
voucher households, compared with households without vouchers (or no
change), would have refuted it. Random assignment makes the two groups
alike in everything except the price, so the difference in consumption is
the effect of the price; households that chose to take up a subsidy might
differ in poverty or tastes.
\end{solution}

\begin{solution}{pb:b1:income-and-substitution-effects:levy}
\textbf{1.} $f = 0.10 \times 20000/2 = 1000$ litres; $o = 0.9 \times 20000 = 18000$
dalens.

\textbf{2.} $f = 0.10 \times 20000/2.5 = 800$ litres.

\textbf{3.} $m^c = 20000 + 0.5 \times 1000 = 20500$; $f = 0.10 \times 20500/2.5 = 820$.

\textbf{4.} \omterm{def:b1:income-and-substitution-effects:effects}{Substitution effect} $820 - 1000 = -180$; \omterm{def:b1:income-and-substitution-effects:effects}{income effect} $800 - 820 = -20$;
together $-200 = 800 - 1000$.

\textbf{5.} $v = m\,a^a(1 - a)^{1-a}p_f^{-a}$, so utility is unchanged when $m$ rises
by the factor $(p_f'/p_f)^a = 1.25^{0.1}$. That gives
$20000 \times 1.022565 = 20451.3$ and $f = 0.10 \times 20451.3/2.5 = 818.1$. Hicks
\omterm{def:b1:income-and-substitution-effects:effects}{substitution effect} $818.1 - 1000 = -181.9$, \omterm{def:b1:income-and-substitution-effects:effects}{income effect} $800 - 818.1 = -18.1$.

\textbf{6.} Fuel takes only $10\%$ of A's budget, so the price rise costs her
little income and the \omterm{def:b1:income-and-substitution-effects:effects}{income effect} is small; the two compensations, 500
and 451.3, differ only by a second-order term.

\textbf{7.} From $v(p, m) = m\,a^a(1 - a)^{1-a}p_f^{-a}$, solving $v = u$ for $m$
gives $e(p, u) = u\,p_f^{a}/\bigl(a^a(1 - a)^{1-a}\bigr)$, and substituting
$u = v(p, m)$ returns $m$.

\textbf{8.} The compensation is $e(p', u_0) - m = m(1.25^{a} - 1)$, since
$e(p, u_0) = m$ and $e$ is proportional to $p_f^a$. A:
$20000(1.25^{0.10} - 1) = 451.3$; B: $30000(1.25^{0.08} - 1) = 540.4$;
C: $50000(1.25^{0.06} - 1) = 673.9$ dalens.

\textbf{9.} $0.5 \times 1000 = 500$, $0.5 \times 1200 = 600$ and $0.5 \times 1500 = 750$ (B uses
$0.08 \times 30000/2 = 1200$ litres and C $0.06 \times 50000/2 = 1500$). The
first-order compensation pays for the old bundle at the new price; a
household can reach its old utility more cheaply by buying less fuel, so
the exact compensation is smaller ($e$ is concave in prices).

\textbf{10.} $a\,e/p_f = 0.10 \times 20451.3/2.5 = 818.1$, A's compensated fuel use in
question 5.

\textbf{11.} After the levy A uses 800, B $0.08 \times 30000/2.5 = 960$ and C
$0.06 \times 50000/2.5 = 1200$ litres: 2,960 litres and a revenue of
$0.5 \times 2960 = 1480$ dalens. The exact compensations add up to
$451.3 + 540.4 + 673.9 = 1665.6$.

\textbf{12.} The levy costs the households $1665.6 - 1480 = 185.6$ dalens more,
in money they would accept instead, than it raises: the deadweight loss of
the levy, which comes from the litres no longer bought.

\textbf{13.} Before: $1000 + 1200 + 1500 = 3700$ litres. After, without the cheque:
2,960, a fall of $740/3700 = 20.0\%$.

\textbf{14.} $f_h = a_h(m_h + c)/2.5$, so
$F = (2000 + 2400 + 3000 + 0.24c)/2.5 = (7400 + 0.24c)/2.5$, with $3c = 0.5F$.

\textbf{15.} $c = F/6$, so $2.5F = 7400 + 0.04F$ and $F = 7400/2.46 = 3008.13$ litres;
$c = 3008.13/6 = 501.36$, about 501.4 dalens.

\textbf{16.} $1 - 3008.13/3700 = 18.7\%$.

\textbf{17.} A $818.1$, B $0.08 \times 30540.4/2.5 = 977.3$, C $0.06 \times 50673.9/2.5 = 1216.2$:
3,011.6 litres, a fall of $1 - 3011.6/3700 = 18.6\%$, almost the same as with the
cheque.

\textbf{18.} The cheque removes the \omterm{def:b1:income-and-substitution-effects:effects}{income effect}, $-s_j\eta_i$ in elasticities,
which is small because fuel takes $6$ to $10\%$ of each budget. The
\omterm{def:b1:income-and-substitution-effects:effects}{substitution effect}, the larger part, stays: the fall goes only from
$20.0\%$ to $18.7\%$.

\textbf{19.} A: $501.4 - 451.3 = 50.1$ dalens better off. B: $501.4 - 540.4 = -39.0$,
worse off. C: $501.4 - 673.9 = -172.5$, worse off.

\textbf{20.} A uses the least fuel and has the highest share. Compare B and C:
B has the higher share but C, which uses more litres, loses more. What
decides is the fuel bought, since the cost of the levy is about the levy
times the litres, while the cheque is the same for all.

\textbf{21.} $50.1 - 39.0 - 172.5 = -161.4$: the three cheques, $3 \times 501.4 = 1504.2$,
fall short of the exact compensations, 1665.6, by the deadweight loss of the
levy when the money comes back.

\textbf{22.} Only for households whose exact compensation is below the cheque,
roughly those that buy less fuel than average: here only A.

\textbf{23.} With $c_h = k\,m_h$ and $k \times 100000 = 0.5F$,
$2.5F = 7400(1 + k) = 7400 + 0.037F$, so $F = 7400/2.463 = 3004.47$ and the fall is
$18.8\%$. The cheque's form barely matters for fuel use, because \omterm{def:b1:income-and-substitution-effects:effects}{income effects}
are small; it matters for who gains.

\textbf{24.} With the cheque, fuel use still falls by $18.7\%$, against $20.0\%$
without it.
\end{solution}
